\subsection{Fire på stribe}
Langt de fleste kender ``fire på stribe" som et kedeligt og simpelt spil, der hører til i børnehaven, men med ganske små modifikationer bliver det dog udfordrende og ikke-trivielt. Vi vil i det følgende vise et par varianter af spillet. Det forudsættes at læseren er bekendt med de sædvanlige regler:
\begin{itemize}
    \item man skiftes til at placere en brik i hver sin farve,
    \item spillepladen er et $6\times7$-bræt.
\end{itemize}

\begin{minipage}{0.45\textwidth}
\centering
\begin{tabular}{|c|c|c|c|c|c|c|}
    \hline
    & & & & & & \\ 
    \hline
    & & & & & & \\
    \hline
    {\scriptsize$\bigcirc$} & & $\times$ & & & & $\oplus$ \\
    \hline
    $\oplus$ & $\times$ & {\scriptsize$\bigcirc$} & & & & {\scriptsize$\bigcirc$} \\
    \hline
    $\times$ & $\oplus$ & $\times$ & $\times$ & {\scriptsize$\bigcirc$} & & $\times$ \\
    \hline
    $\times$ & $\times$ & $\oplus$ & $\times$ & $\times $& {\scriptsize$\bigcirc$} & {\scriptsize$\bigcirc$} \\
    \hline
\end{tabular}

{\scriptsize$\bigcirc$} vinder i \textit{torus} varianten.
\end{minipage}
\begin{minipage}{0.45\textwidth}
\centering
\begin{tabular}{|c|c|c|c|c|c|c|}
    \hline
    & & & & & & \\ 
    \hline
    & & & & & & \\
    \hline
    {\scriptsize$\bigcirc$} & & $\times$ & & & & {\scriptsize$\bigcirc$} \\
    \hline
    $\oplus$ & $\times$ & {\scriptsize$\bigcirc$} & & & & $\oplus$ \\
    \hline
    $\times$ & $\oplus$ & $\times$ & $\times$ & {\scriptsize$\bigcirc$} & & $\times$ \\
    \hline
    $\times$ & $\times$ & $\oplus$ & $\times$ & $\times $& {\scriptsize$\bigcirc$} & {\scriptsize$\bigcirc$} \\
    \hline
\end{tabular}

{\scriptsize$\bigcirc$} vinder i \textit{Möbius} varianten.
\end{minipage}
\newline\newline


\begin{minipage}{0.45\textwidth}
\centering
\begin{tabular}{|c|c|c|c|c|c|c|}
    \hline
    & & & & & & \\ 
    \hline
    & & & & & & \\
    \hline
    {\scriptsize$\bigcirc$} & & $\otimes$ & & & & {\scriptsize$\bigcirc$} \\
    \hline
    {\scriptsize$\bigcirc$} & $\otimes$ & {\scriptsize$\bigcirc$} & & & & {\scriptsize$\bigcirc$} \\
    \hline
    $\otimes$ & {\scriptsize$\bigcirc$} & $\times$ & $\times$ & {\scriptsize$\bigcirc$} & & $\times$ \\
    \hline
    $\times$ & $\otimes$ & {\scriptsize$\bigcirc$} & $\times$ & $\times $& {\scriptsize$\bigcirc$} & {\scriptsize$\bigcirc$} \\
    \hline
\end{tabular}

$\times$ vinder i \textit{bounce} varianten.
\end{minipage}
\begin{minipage}{0.45\textwidth}
\centering
\begin{tabular}{|c|c|c|c|c|c|c|}
    \hline
    & & & & & & \\ 
    \hline
    & & & & & & \\
    \hline
    {\scriptsize$\bigcirc$} & & $\times$ & & & & {\scriptsize$\bigcirc$} \\
    \hline
    {\scriptsize$\bigcirc$} & $\otimes$ & {\scriptsize$\bigcirc$} & & & & {\scriptsize$\bigcirc$} \\
    \hline
    $\otimes$ & {\scriptsize$\bigcirc$} & $\otimes$ & $\times$ & {\scriptsize$\bigcirc$} & & $\times$ \\
    \hline
    $\times$ & $\otimes$ & {\scriptsize$\bigcirc$} & $\times$ & $\times $& {\scriptsize$\bigcirc$} & {\scriptsize$\bigcirc$} \\
    \hline
\end{tabular}

$\times$ vinder i \textit{square} varianten.
\end{minipage}

\noindent\textbf{Torus}\newline
I \textit{torus}, skal man forestille sig at siderne er klistret sammen, og ligeledes med top og bund. Geometrisk spiller man nu på en torus. Det gør, at man kan ``gå ud af den ene side og komme ind ad den anden", hvilket blandt andet giver vindermulighederne illustreret som den venstre øverste opstilling, hvor $\oplus$ indikere den vindende stribe for $\scriptsize\bigcirc$.\newline

%Hvis man har svært ved at overskue spillet, kan man forestille sig, at man tager kopier af den originale spilleplade, og placerer disse ovenover, nedenunder, venstre og højre for. Bemærk dog, at de diagonale hjørner hænger sammen.

% Da tyngdekraften er en essentiel del af spillet, er det naturligvis en forudsætning for de sidste tre muligheder at der er fyldt op med brikker under de illustrerede vinderbrikker.

\noindent\textbf{Möbius}\newline
En lignende variant er \textit{Möbius}, hvor top og bund er limet sammen men siderne er ``vredet" og klistret sammen. Så hjørnet i venstre bund er klistret på højre top, og hjørnet i venstre top er klistret på højre bund. Dette bliver hurtigt en ``hjerne-vrider" når man skal finde ud af om man har vundet. En vinderstribe i Möbius er vist øverst til højre i illustrationen.\newline

\noindent\textbf{Bounce}\newline
En tredje variant er \textit{bounce}, hvor man kan bande på væggene, hvilket i det tidligere eksempel ville betyde at $\times$ havde vundet i stedet for $ $, som kan nederst til venstre på illustrationen, hvor $\otimes$ indikere vinderstriben for $\times$.\newline

\noindent\textbf{Square}\newline
Endnu en variant er \textit{square}, hvor det ikke handler om at få fire på stribe, men derimod fire brikker til at udgøre hjørnerne i et (ikke nødvendigvis akseparallelt) kvadrat, som i den nederstehøjre opstilling i illustrationen. Dette kan f.eks. kombineres med torus eller Möbius, hvorved man kan opnå nogle aldeles besynderlige vinderkonstellationer. Det er naturligvis et krav, at der altid indgår fire forskellige brikker i en vinderkonstellation.\newline

\noindent\textbf{Hold varianter}\newline
Når man behersker alle disse varianter, er det på tide man spiller i hold. En typisk sammensætning er to mod to, hvor det er forbudt at snakke om spillet indbyrdes. Det betyder, at man bliver nødt til at håbe ens medspiller har forståets ens taktik, inden modstanderne gætter denne. Man skiftes selvfølgelig indbyrdes til at droppe brikker i rammen, ligesom holdene skiftes.
 
Alternativt, kan man, i tilfælde af et ulige antal spillere, skiftes til at droppe en brik, mens farveordenen opretholdes, således at farverne skiftes til at komme i. Herved vil hver spiller skiftevis droppe forskellige farver i rammen. Den person, som lægger vinderbrikken, vinder, resten taber. God fornøjelse.
